\scanpage{37}
Soit\deleted{\unsure{ent|es}} $S$ schéma, \deleted{$u\in M_2(S)=M_2(\Gamma(S,\mathcal O_S))$, supposons} \inserted{$V$ \deleted{fibré} \inserted{Module} vectoriel de rang 2, $u\in\operatorname{End}(V)$}
\[
 \forall s\in S,\qquad u_s\notin k(s).1 \tag*{(1)\deleted{\unsure{'}}}
\]
Considérons le commutant de $u$ dans $\underline{\operatorname{End}}(V)$ (le foncteur). Je dis qu'il est égal au sous-fibré vectoriel $\mathcal O_S\,\mathrm{Id}+\mathcal O_S\,u$ de $\underline{\operatorname{End}}(V)$ -- donc le commutant dans $\operatorname{GL}(V)$ est formé des $\alpha\,\mathrm{id}+\beta u$, avec
\[
 \det(\alpha\,\mathrm{id}+\beta u)=\alpha^2+\alpha\beta\operatorname{Tr}u+\beta^2\det u\ \deleted{\unsure{\ne}}\quad\text{inv.} \tag*{(2)}
\]
donc le commutant (ou stabilisateur) dans $\underline{\operatorname{GP}}(V)$ s'identifie au
\[
 C_{\underline{\operatorname{GP}}(V)}(u)\simeq\mathbb P^1-Q_{\operatorname{Tr}u,\det u}
\]
où $Q_{\deleted{\unsure{\operatorname{Tr}u}}\tau,\delta}$ désigne, de façon \unsure{général|générale}, la ``\unsure{0-|o-}quadrique'' de $\mathbb P^1$ définie par l'équation quadratique en $\alpha,\beta$
\[
 \alpha^2+\tau.\alpha\beta+\delta.\beta^2\deleted{\det u}
\]
En l'occurrence, c'est un diviseur de degré 2 \deleted{\unsure{donc}} qui est étale sss
\[
 \tau^2-4\delta\qquad\text{inv.}
\]
ce qui, dans le cas $\tau=\operatorname{Tr}u\deleted{\unsure{=0}},\delta=\det u$, signifie que les val. propres de $u$ fibre par fibre sont distinctes,
% archive sheet number: 36
